\documentclass[12pt,a4paper]{article}
\usepackage[margin=25mm, headheight=15pt, headsep=10pt]{geometry}
\usepackage{fontspec}
\usepackage[table]{xcolor} % Note: table option is required for rowcolors
\usepackage{titlesec}
\usepackage{tocloft}
\usepackage{fancyhdr}
\usepackage{booktabs}
\usepackage{tcolorbox}
\tcbuselibrary{skins, breakable}
\usepackage[colorlinks=true, linkcolor=emerald, urlcolor=emerald, bookmarks=true]{hyperref}

\definecolor{emerald}{RGB}{0,102,51}
\definecolor{darkred}{RGB}{139,0,0}
\definecolor{lightgray}{RGB}{245,245,245}

\usepackage[nil]{babel}
\babelprovide[import, main]{bengali}
\babelprovide[import]{arabic}
\babelfont{rm}[Renderer=HarfBuzz,Scale=1.0]{Noto Serif Bengali}
\babelfont{sf}[Renderer=HarfBuzz]{Noto Sans Bengali}
\babelfont{tt}[Renderer=HarfBuzz]{Noto Sans Bengali}
\babelfont[arabic]{rm}[Renderer=HarfBuzz,Scale=1.1]{Noto Naskh Arabic}

% Section styling
\titleformat{\section}{\Large\bfseries\color{emerald}}{\thesection.}{0.5em}{}
\titleformat{\subsection}{\large\bfseries\color{emerald!85!black}}{\thesubsection.}{0.5em}{}
\titleformat{\subsubsection}{\normalsize\bfseries\color{emerald!70!black}}{\thesubsubsection.}{0.5em}{}
\titlespacing*{\section}{0pt}{18pt}{8pt}

% Fancy Header/Footer setup
\pagestyle{fancy}
\fancyhf{} % clear all header and footer fields
\fancyhead[L]{\color{emerald}\small\textbf{\nouppercase{\leftmark}}}
\fancyfoot[L]{\scriptsize\color{black!50}{\lat{AI}}-সংকলিত, আলেম-যাচাইকৃত নয়}
\fancyfoot[C]{\color{emerald!70!black}\thepage}
\renewcommand{\headrulewidth}{0.4pt}
\renewcommand{\headrule}{\hbox to\headwidth{\color{emerald}\leaders\hrule height \headrulewidth\hfill}}
\renewcommand{\footrulewidth}{0pt}

% Custom boxes
% 1. Ayat/Hadith Box
\newtcolorbox{ayatbox}[1][]{
    enhanced, breakable, 
    colback=emerald!5, 
    colframe=emerald, 
    boxrule=0.5pt, 
    arc=3pt, 
    left=10pt, right=10pt, top=10pt, bottom=10pt,
    #1
}

% 2. Quote/Translation Box
\newtcolorbox{quotebox}[1][]{
    enhanced, breakable,
    frame hidden,
    colback=lightgray,
    borderline west={3pt}{0pt}{emerald},
    left=15pt, right=10pt, top=10pt, bottom=10pt,
    sharp corners,
    #1
}

% 3. AI-generated notice box
\newtcolorbox{warnbox}[1][]{
    enhanced, breakable,
    colback=red!4,
    colframe=darkred,
    boxrule=0.8pt,
    arc=3pt,
    left=10pt, right=10pt, top=10pt, bottom=10pt,
    #1
}

% Commands
\newcommand{\arab}[1]{\foreignlanguage{arabic}{#1}}
\newcommand{\kib}[1]{\textcolor{emerald}{\textbf{#1}}}

% Latin (English) fallback: Noto Serif Bengali has NO Latin glyphs
\newfontfamily\latfont[Renderer=HarfBuzz]{Noto Serif}
\newcommand{\lat}[1]{{\latfont #1}}

\setlength{\parskip}{5pt}
\setlength{\parindent}{0pt}

% Fix for missing bullet in Noto Serif Bengali
\renewcommand{\labelitemi}{$\bullet$}



\begin{document}
\begin{center}{\Large\bfseries\color{emerald} পার্ট ৫গ — মাসআলা: সফর (কসর-জমা), অসুস্থ, বসে-শুয়ে নামাজ, নারীদের সালাত}\end{center}
\vspace{1em}
\section*{পার্ট ৫গ — মাসআলা: সফর (কসর-জমা), অসুস্থ, বসে-শুয়ে নামাজ, নারীদের সালাত}

\begin{warnbox}
 \textbf{গুরুত্বপূর্ণ নোটিশ (\lat{Important} \lat{Notice}):} এই কন্টেন্টটি \lat{AI} (কৃত্রিম বুদ্ধিমত্তা) দিয়ে প্রস্তুত ও সংকলিত। \lat{AI} কঠোর সূত্র-মান নীতি মেনে শুধু নির্ভরযোগ্য উৎস থেকে তথ্য সংগ্রহ করেছে — কেবল সহীহ/হাসান হাদিস ও সহীহ সনদের আসার রাখা হয়েছে, দুর্বল সনদ বাদ দেওয়া হয়েছে। তবে এটি কোনো আলেম বা মুহাদ্দিস কর্তৃক যাচাইকৃত নয়।
\end{warnbox}

\begin{quote}
\textbf{সম্পর্কিত ফাইল:} পার্ট ২ (পূর্বশর্ত), পার্ট ৩ (শ্রেণিবিভাগ), পার্ট ৫ক (জামাত), পার্ট ৫খ (সাহু সিজদা), পার্ট ৬ (মাযহাব-তুলনা)।
\end{quote}

\begin{quote}
\textbf{যে প্রশ্নের উত্তর এই অংশে:} সফরে কীভাবে নামাজ পড়ব (কসর-জমা)? সফরের দূরত্ব কত? অসুস্থ হলে বসে/শুয়ে নামাজের নিয়ম? নারীর নামাজের বিশেষ মাসআলা (হায়েয-নিফাস, নারী-জামাত)?
\end{quote}

\medskip\hrule\medskip

\subsection*{১. সফরের নামাজ — কসর ও জমা}

\subsubsection*{১.১ কসর (সংক্ষিপ্ত) — ভিত্তি}

\textbf{সূরা আন-নিসা ৪:১০১:}

\begin{center}\arab{وَإِذَا ضَرَبْتُمْ فِى ٱلْأَرْضِ فَلَيْسَ عَلَيْكُمْ جُنَاحٌ أَن تَقْصُرُوا۟ مِنَ ٱلصَّلَوٰةِ إِنْ خِفْتُمْ أَن يَفْتِنَكُمُ ٱلَّذِينَ كَفَرُوٓا۟ ۚ إِنَّ ٱلْكَـٰفِرِينَ كَانُوا۟ لَكُمْ عَدُوًّا مُّبِينًا}\end{center}

\textbf{উচ্চারণ:} ওয়া ইযা- দ্বারাবতুম ফিল-আরদ্বি ফালাইসা 'আলাইকুম জুনাহুন আন তাক্বসুরূ মিনাস-সলাতি ইন খিফতুম আইয়্যাফতিনাকুমুল্লাযীনা কাফারূ। ইন্নাল-কাফিরীনা কা-নূ লাকুম 'আদুউয়্যাম মুবীনান।

\textbf{বাংলা অর্থ (\lat{pure}):}
\textbf{\lat{“}আর যখন তোমরা পৃথিবীতে ভ্রমণ করো, তখন নামাজ সংক্ষেপ করায় তোমাদের কোনো গুনাহ নেই — যদি ভয় করো যে কাফেররা তোমাদের ফিতনায় ফেলবে। নিশ্চয়ই কাফেররা তোমাদের প্রকাশ্য শত্রু।\lat{”}}

\textbf{ব্যাখ্যা:} আয়াতে কসরের শর্ত বলা হয়েছে \lat{“}ভয়\lat{”} — কিন্তু নবীজির (সা.) আমলে দেখা যায়, শান্তির যুগেও তিনি সফরে কসর করেছেন। তাই ফকিহরা মনে করেন — \textbf{সফরে কসর সুন্নাতে মুআক্কাদা} (হানাফি: ওয়াজিব নয়, তবে অত্যন্ত গুরুত্বপূর্ণ); জমহুর: রুখসত (অনুমতি) — তবে উত্তম হলো কসর।

\textbf{সহীহ হাদিস (সহীহ বুখারি ১০৮২; সহীহ মুসলিম ৬৮৬):}

\begin{quote}
আয়েশা (রাঃ) থেকে: \lat{“}সফরের নামাজ দুই রাকাত করে\lat{”} — যখন নবীজি (সা.) সফরে থাকতেন, চার রাকাত-ফরজ দুই রাকাত পড়তেন (যোহর-আসর-ইশা), ফজর-মাগরিব অপরিবর্তিত।
\end{quote}

\begin{quote}
\textbf{কসরের নিয়ম:} যোহর, আসর, ইশা — \textbf{৪ রাকাতের পরিবর্তে ২ রাকাত}। ফজর (২) ও মাগরিব (৩) কসর হয় না। সুন্নত রাকাতও সফরে ছাড়া/কসর করা যায় (মাযহাবভেদে; ফজরের ২ রাকাত সুন্নত নবীজি সফরেও পড়েছেন — সহীহ বুখারি ১১৬৫)।
\end{quote}

\subsubsection*{১.২ সফরের দূরত্ব — মাযহাবগত অবস্থান}

\begin{table}[h]\centering\small
\begin{tabular}{ll}
মাযহাব & কসরযোগ্য সফরের ন্যূনতম দূরত্ব \\
\hline
\textbf{হানাফি} & ৩ দিনের (৭৮ কিমি-প্রায়, হাটার/সামান্য বাহনের) সফর — প্রসিদ্ধ: ৩ দিন-২ রাত / ১২০.৫ কিমি (অনেক হিসাবে) \\
\textbf{শাফেয়ি} & ৪৮ কিমি-প্রায় (বারিদ=২৪ মাইল) \\
\textbf{মালেকি} & ৪৮ মাইল-প্রায় (৮০ কিমি) \\
\textbf{হাম্বলি} & ৮০ কিমি-প্রায় (একটি রাতের হাঁটা সফর) \\
\hline
\end{tabular}
\end{table}

\begin{quote}
\textbf{সারমর্ম:} সব মাযহাবের মিল — \textbf{সাধারণ সফরের দূরত্ব} (এক-দুই দিনের ভ্রমণ) হলে কসর জায়েয; সীমা সংখ্যায় আলাদা। আধুনিক যান-বাহনে দূরত্ব হিসাব করা হয় — নয়তো \textbf{"সফর" শব্দের সাধারণ প্রচলিত অর্থ} অনুযায়ী (নির্বাচিত ফতোয়ার ভিত্তি)।
\end{quote}

\subsubsection*{১.৩ জমা (একত্র করা) — যোহর-আসর, মাগরিব-ইশা}

\textbf{সহীহ হাদিস (সহীহ মুসলিম ৭০৫-৭০৬):}

\begin{quote}
নবীজি (সা.) সফরে যোহর ও আসর — এবং মাগরিব ও ইশা — \textbf{একই সময়ে (জমা)} পড়েছেন (তাবুক সফর, এবং অন্যান্য সময়)।
\end{quote}

\begin{itemize}
\item \textbf{জমা দুই ধরনের:}
\item \textbf{জমা তাকদিম:} প্রথম নামাজের সময়ে দুইটা (যেমন যোহরের সময়ে যোহর+আসর)।
\item \textbf{জমা তাখির:} দ্বিতীয় নামাজের সময়ে (যেমন ইশার সময়ে মাগরিব+ইশা)।
\item \textbf{জমা বৈধ:} সফরে, এবং জমহুর মতে প্রবল বৃষ্টি/আপদেও (মালেকি আরও বিস্তৃত)। হানাফিতে \textbf{সফরে জমা জায়েয} (সহীহ হাদিসে প্রমাণিত), তবে ইমামদের মধ্যে \textbf{আসরের সাথে যোহর জমা করার বৈধতায়} মতভেদ আছে (বিস্তারিত পার্ট ৬)।
\item \textbf{জমা করার ক্রম:} তাকদিমে প্রথমে প্রথমটি, তাখিরে প্রথমটি-পরে দ্বিতীয়টি — ক্রম বজায় রাখতে হবে।
\end{itemize}
\subsubsection*{১.৪ সফরের মাসআলা-সংক্ষেপ}

\begin{itemize}
\item মুসাফির কসর করবেন — \textbf{সফরের পুরো সময়} (মুকিম ঘোষণা না হওয়া পর্যন্ত)।
\item \textbf{মুসাফির-ইমাম ও মুকিম-মুক্তাদি:} ইমাম কসর করলে মুক্তাদিও কসর-সংক্ষেপে অনুসরণ (জামাতের রাকাত সংখ্যা ইমামের অনুযায়ী) — কিন্তু ইমাম সালামের পর মুকিম \textbf{বাকি রাকাত পূর্ণ করবে} (মাসবুক-নিয়মে)।
\item \textbf{সফরে জুমা:} জুমা ফরজ নয়; মুসাফির জুমার বদলে জোহর পড়ে (বিস্তারিত পার্ট ৫ঘ)।
\end{itemize}
\medskip\hrule\medskip

\subsection*{২. অসুস্থ ব্যক্তির নামাজ}

\subsubsection*{২.১ বসে/শুয়ে নামাজ — ভিত্তি}

\textbf{সহীহ হাদিস (সহীহ বুখারি ১১১৭):}

\begin{quote}
ইমরান ইবনে হুসাইন (রাঃ) থেকে: আমার ছিল অর্শ (ববাসির) রোগ। আমি নবীজি (সা.)-কে বললাম: আমি বসে নামাজ পড়ি, এর হুকুম কী? নবীজি (সা.) বললেন: \textbf{\lat{“}দাঁড়িয়ে পড়ো; না পারলে বসে; না পারলে কাত হয়ে।\lat{”}}
\end{quote}

\begin{quote}
\textbf{মূলনীতি:} কুদরত অনুযায়ী — দাঁড়ানো অসম্ভব হলে বসে, বসাও অসম্ভব হলে শুয়ে (পার্শ্বে কাত হয়ে / সিজদা ইঙ্গিতে)। সিজদা নিচু করা সম্ভব না হলে \textbf{ইশারা} (সংকেতে) — সিজদা রুকুর চেয়ে বেশি নিচু করে ইশারা।
\end{quote}

\textbf{সহীহ হাদিস (সহীহ বুখারি ১১১৮):}

\begin{quote}
নবীজি (সা.) অসুস্থ অবস্থায় বসে নামাজ পড়েছেন; তাঁর পেছনে সাহাবিরা দাঁড়িয়ে পড়েন — নবীজি (সা.) ইশারায় বললেন বসে পড়ো।
\end{quote}

\subsubsection*{২.২ অসুস্থ ব্যক্তির নিয়ম-সংক্ষেপ}

\begin{enumerate}
\item \textbf{কিরাআত-রুকু-সিজদা:} সামর্থ্য অনুযায়ী — দাঁড়িয়ে, বসে, বা শুয়ে; সিজদা ইঙ্গিতে।
\item \textbf{বসে রুকু-সিজদা:} বসে থাকলে রুকু সামান্য ঝুঁকে, সিজদা আরও ঝুঁকে (নিচুতে)।
\item \textbf{সিজদা-ইশারা:} সিজদা দিতে না পারলে কপালে হাত-মাথা নিচু করে ইশারা; রুকুর ইশারার চেয়ে সিজদা নিচু।
\item \textbf{ফরজ নামাজ রাহাত দেয় না} — ওয়াক্তে প্রতিটি ফরজ অবশ্যই আদায় (মুমিনদের উপর কিতাব-মাওকুতা) — অসুস্থ অবস্থায়ও ওয়াক্ত পেরিয়ে যাওয়া থেকে রেহাই পেতে যথাসাধ্য।
\end{enumerate}
\medskip\hrule\medskip

\subsection*{৩. নারীর নামাজের বিশেষ মাসআলা}

\subsubsection*{৩.১ হায়েয (মাসিক) ও নিফাস (সন্তান-পরবর্তী রক্ত)}

\textbf{সহীহ হাদিস (সহীহ বুখারি ৩০৪):}

\begin{quote}
আয়েশা (রাঃ) থেকে: নবীজি (সা.) বলেছেন — \textbf{\lat{“}কি! তুমি কি নামাজ পড়ো না?\lat{”}} (হায়েয অবস্থায়)। আয়েশা বললেন — নবীজি (সা.)-কে জিজ্ঞেস করলেন। তিনি বললেন — \textbf{\lat{“}এটা (হায়েয) আল্লাহর নির্ধারিত নিয়ম (লেখা), যা আদম-কন্যাদের উপর; তাই হায়েয অবস্থায় নামাজ ছাড়ো, পরে কাযা কোরো না (নামাজের কাযা নেই; রোজার কাযা আছে)।\lat{”}}
\end{quote}

\begin{itemize}
\item \textbf{নিয়ম:} হায়েয-নিফাস অবস্থায় \textbf{নামাজ, রোজা, কুরআন-স্পর্শ} (মাযহাবগত) থেকে বিরত; নামাজের \textbf{কাযা নেই}, রোজার কাযা আছে (সহীহ মুসলিম ৩৩৫-এ আয়েশা থেকে স্পষ্ট)।
\item \textbf{হায়েয শেষে:} গোসল করে স্বাভাবিক ইবাদতে ফেরা।
\item \textbf{বিস্তারিত মাসআলা} (হায়েযের ন্যূনতম-সর্বোচ্চ দিন, ইস্তিহাজা) — পার্ট ৬-এ মাযহাব-তুলনায়।
\end{itemize}
\subsubsection*{৩.২ নারী-জামাত}

\begin{itemize}
\item \textbf{নারীর ইমামতি:} পুরুষের ইমামতিতে নারী \textbf{জায়েয নয়} (সব মাযহাবের জমহুর)। নারীদের নিজেদের মাঝে — \textbf{নারী-ইমামতি জায়েয} (মাযহাবগত প্রসিদ্ধ: উম্মে ওয়ারাকা (রাঃ) তাঁর গৃহের নারীদের ইমামতি করেছেন — সুনানে আবু দাউদ ৫৯১ — সহীহ সনদে এই ঘটনা আছে)।
\item \textbf{নারী মসজিদে:} ফরজ জামাতে মসজিদে নারীর অংশগ্রহণ জায়েয ও ফজিলতপূর্ণ (বুখারি ৯০০; পার্ট ৫ক)। নবীজি-যুগে নারীরা জামাতে অংশ নিতেন (উম্মে আতিয়্যাহর ঈদ-ঘটনা — পার্ট ৮)।
\item \textbf{কাতার:} নারীরা পেছনের কাতারে; পুরুষের সাথে মিশ্র কাতার বৈধ নয়।
\end{itemize}
\subsubsection*{৩.৩ নারীর সতর-পোশাক}

\begin{itemize}
\item মুখ ও দুই হাত ছাড়া সব ঢেকে নামাজ (আরও বিশদ পার্ট ২-এর সতর অংশ; মাযহাবগত সামান্য পার্থক্য)।
\end{itemize}
\medskip\hrule\medskip

\subsection*{৪. তায়াম্মুম-মাসেহ-এর বিশেষ অবস্থা (সফর-অসুস্থ প্রসঙ্গে)}

\begin{itemize}
\item \textbf{তায়াম্মুম:} পানি না পেলে/অক্ষম হলে — পবিত্র মাটি দিয়ে চেহারা ও দুই হাত মাসেহ (সূরা ৪:৪৩; বুখারি ৩৩৫, ৩৪৮ — পার্ট ২)। \textbf{প্রতিটি ফরজের জন্য} নতুন তায়াম্মুম (মাযহাবগত: হানাফি এক নামাজে এক তায়াম্মুম যথেষ্ট; জমহুর প্রতিটি নামাজে নতুন)।
\item \textbf{মাসেহ খুফফাইন:} মোজার উপর ২৪ ঘণ্টা (মুকিম) / ৭২ ঘণ্টা (মুসাফির) — (সহীহ মুসলিম ২৭৬; পার্ট ২)।
\end{itemize}
\medskip\hrule\medskip

\subsection*{৫. প্রশ্নোত্তর}

\begin{table}[h]\centering\small
\begin{tabular}{ll}
প্রশ্ন & উত্তর \\
\hline
সফরে ৪ রাকাত কসর করা কি ওয়াজিব? & সুন্নাতে মুআক্কাদা (হানাফি)/রুখসত (জমহুর) — তবে কসর করাই নবীজির আমল \\
সফরের দূরত্ব কত? & মাযহাবভেদে ৪৮–১২০ কিমি-মধ্য; সাধারণ ভ্রমণ-দূরত্ব হলে কসর জায়েয \\
জমা কখন? & সফরে, ও প্রবল বৃষ্টি-আপদে (জমহুর); হানাফিতে সফরে জমা জায়েয \\
অসুস্থ বসে পড়তে পারি? & পারি — দাঁড়াতে না পারলে (বুখারি ১১১৭) \\
হায়েয অবস্থায় নামাজ? & ছাড়বে, কাযা নেই, রোজা-কাযা আছে (বুখারি ৩০৪) \\
নারী পুরুষদের ইমামতি? & জায়েয নয় (জমহুর) — নারী নারীদের ইমামতি জায়েয \\
\hline
\end{tabular}
\end{table}

\medskip\hrule\medskip

\subsection*{৬. রেফারেন্স}

\begin{itemize}
\item কুরআন: আন-নিসা ৪:১০১, ৪:৪৩; আল-মায়িদাহ ৫:৬
\item সহীহ বুখারি — ৩০৪, ৩৩৫, ৩৪৮, ১০৮২, ১১১৭, ১১১৮, ১১৬৫, ৯০০
\item সহীহ মুসলিম — ২৭৬, ৩৩৫, ৬৮৬, ৭০৫-৭০৬
\item সুনানে আবু দাউদ — ৫৯১ (উম্মে ওয়ারাকা), ৬৬০
\item আল-ফিকহুল ইসলামি ওয়া আদিল্লাতুহু (আল-যুহাইলি) — মাযহাব-সারণি
\end{itemize}
\begin{quote}
\textbf{দ্রষ্টব্য:} জুমা-ঈদ-জানাজা-নফল (পার্ট ৫ঘ) ও বিতর্কিত মাসআলার মাযহাব-তুলনা (পার্ট ৬) আলাদা ফাইল; হায়েয-নিফাসের সম্পূর্ণ ফিকহি বিশদ ও নারী-সংক্রান্ত বিতর্কগুলো পার্ট ৬-এ।
\end{quote}

\end{document}
